\documentclass[11pt,a4paper]{article}
\usepackage[utf8]{inputenc}
\usepackage[vietnamese]{babel}
\usepackage{amsmath,amssymb}
\usepackage{geometry}
\geometry{a4paper, total={180mm,260mm}, left=15mm, right=15mm, top=18mm, bottom=18mm}
\usepackage{tikz}
\usetikzlibrary{shapes.geometric, arrows.meta, positioning, calc}
\usepackage{tcolorbox}
\tcbuselibrary{skins,breakable}
\usepackage{xcolor}
\usepackage{multicol}
\usepackage{array}
\usepackage{fontawesome}
\usepackage{fancyhdr}

\pagestyle{fancy}
\fancyhf{}
\renewcommand{\headrulewidth}{0.5pt}
\renewcommand{\footrulewidth}{0.5pt}
\fancyhead[L]{\small\bfseries\color{teal}\faBook\ SỔ TAY KIẾN THỨC TOÁN 3}
\fancyhead[R]{\small\bfseries\color{teal}Lớp Toán Thầy Lĩnh -- loptoanthaylinh.vn}
\fancyfoot[C]{\small\color{gray}-- Sổ tay kiến thức: Bài 71. Chia số có năm chữ số cho số có một chữ số --}

\definecolor{primary}{HTML}{005F73}
\definecolor{secondary}{HTML}{0A9396}
\definecolor{accent}{HTML}{EE9B00}
\definecolor{danger}{HTML}{AE2012}
\definecolor{boxbg}{HTML}{F4F6F8}

\begin{document}

% --- BANNER TIÊU ĐỀ ---
\begin{tcolorbox}[
  enhanced,
  colback=white,
  colframe=primary,
  arc=3mm,
  boxrule=1.2pt,
  left=4mm, right=4mm, top=2.5mm, bottom=2.5mm,
  drop shadow=black!10!white
]
  \centering
  {\Large\bfseries\color{primary}\faBookmark\ BÀI 71: CHIA SỐ CÓ 5 CHỮ SỐ CHO SỐ CÓ 1 CHỮ SỐ}\\[1.5mm]
  {\small\itshape\color{black!75}Quy trình Chia - Nhân - Trừ từ trái sang phải, hạ từng chữ số và xử lý triệt để chữ số 0 ở thương trên 1 trang A4}
\end{tcolorbox}

\vspace{-2mm}


\noindent
\begin{minipage}[t]{0.49\textwidth}\vspace{0pt}
\begin{tcolorbox}[
  enhanced,
  colback=blue!4!white,
  colframe=blue!75!black,
  arc=2mm, boxrule=0.8pt,
  left=2mm, right=2mm, top=1.8mm, bottom=1.8mm,
  title={\footnotesize\bfseries\color{white}\faCalculator\ 1. Quy trình Đặt tính và Chia}
]
{\footnotesize
\begin{itemize}\setlength{\itemsep}{1.5pt}\setlength{\leftskip}{-2.5mm}
  \item Thực hiện chia theo thứ tự từ \textbf{trái sang phải}:
  $$\text{Chục nghìn} \to \text{Nghìn} \to \text{Trăm} \to \text{Chục} \to \text{Đơn vị}$$
  \item Mỗi lượt chia gồm đủ 3 nhịp: Chia $\to$ Nhân $\to$ Trừ.
  \item \textbf{Lượt chia đầu:} Nếu chữ số hàng chục nghìn nhỏ hơn số chia, lấy cả 2 chữ số đầu để chia.
  \item \textbf{Các lượt chia tiếp theo:} Mỗi lần chỉ hạ đúng $\mathbf{1}$ chữ số.
\end{itemize}
}
\end{tcolorbox}
\end{minipage}
\hfill
\begin{minipage}[t]{0.49\textwidth}\vspace{0pt}
\begin{tcolorbox}[
  enhanced,
  colback=teal!4!white,
  colframe=teal!75!black,
  arc=2mm, boxrule=0.8pt,
  left=2mm, right=2mm, top=1.8mm, bottom=1.8mm,
  title={\footnotesize\bfseries\color{white}\faKey\ 2. Kĩ thuật chữ số 0 ở thương}
]
{\footnotesize
\begin{itemize}\setlength{\itemsep}{1.5pt}\setlength{\leftskip}{-2.5mm}
  \item Khi hạ một chữ số ở bất kì lượt nào sau lượt đầu, nếu số bị chia nhỏ hơn số chia:
  $$\implies \text{Phải viết } \mathbf{0} \text{ vào thương}$$
  rồi mới hạ tiếp chữ số liền kề!
  \item \textbf{Chia số tròn chục nghìn, tròn trăm nghìn nhẩm:}\\
  $80\,000 : 4 = \mathbf{20\,000}$.\\
  $100\,000 : 2 = \mathbf{50\,000}$.
\end{itemize}
}
\end{tcolorbox}
\end{minipage}

\vspace{-2.5mm}
% --- KHỐI 3: SƠ ĐỒ TƯ DUY TIKZ ---
\noindent
\begin{tcolorbox}[
  enhanced,
  colback=orange!4!white,
  colframe=orange!80!black,
  arc=2mm, boxrule=0.8pt,
  left=2mm, right=2mm, top=1.5mm, bottom=1.5mm,
  title={\footnotesize\bfseries\color{white}\faSitemap\ SƠ ĐỒ TƯ DUY: VÒNG LẶP CHIA SỐ CÓ 5 CHỮ SỐ}
]
\centering
\begin{tikzpicture}[
  box/.style={draw, rounded corners=2pt, align=center, font=\scriptsize\bfseries, inner sep=2.5pt, minimum height=0.75cm},
  arr/.style={-{Stealth[length=2mm]}, thick}
]
  \node[box, fill=purple!15, draw=purple!80!black, text width=2.6cm] (d1) at (0, 0) {LƯỢT ĐẦU\\Lấy 1 hoặc 2 số\\Chia được thương đầu};
  \node[box, fill=teal!15, draw=teal!80!black, text width=2.6cm] (d2) at (3.5, 0) {NHÂN \& TRỪ\\Nhân thương với $d$\\Trừ để tìm số dư};
  \node[box, fill=blue!12, draw=blue!80!black, text width=2.6cm] (d3) at (7.0, 0) {HẠ 1 CHỮ SỐ\\Nếu bé hơn số chia\\$\implies$ Viết 0 vào thương};
  \node[box, fill=orange!15, draw=orange!80!black, text width=2.6cm] (d4) at (10.5, 0) {LẶP LẠI\\Đến hàng đơn vị\\Số dư cuối $< d$};

  \draw[arr, purple!80!black] (d1) -- (d2);
  \draw[arr, teal!80!black] (d2) -- (d3);
  \draw[arr, blue!80!black] (d3) -- (d4);
\end{tikzpicture}
\end{tcolorbox}

\vspace{-2.5mm}
% --- KHỐI 4: VÍ DỤ & MẸO TRÁNH BẪY ---
\noindent
\begin{minipage}[t]{0.49\textwidth}\vspace{0pt}
\begin{tcolorbox}[
  enhanced,
  colback=green!4!white,
  colframe=green!60!black,
  arc=2mm, boxrule=0.8pt,
  left=2mm, right=2mm, top=1.8mm, bottom=1.8mm,
  title={\footnotesize\bfseries\color{white}\faCheckCircle\ Ví dụ thực hành chuẩn mực}
]
{\footnotesize
\textbf{1. Phép chia hết có số 0 ở thương:}\\
$63045 : 3 = \mathbf{21015}$.\\
-- $6 : 3 = 2$; hạ $3$, $3 : 3 = 1$.\\
-- Hạ $0$, $0 < 3 \implies$ \textbf{viết 0 vào thương}.\\
-- Hạ $4$, $4 : 3 = 1$ (dư $1$).\\
-- Hạ $5$ được $15$; $15 : 3 = 5$ (dư $0$).\\
\textbf{2. Phép chia có dư:}\\
$45826 : 4 = \mathbf{11456}$ (dư $2$).\\
Thử lại: $11\,456 \times 4 + 2 = 45\,824 + 2 = 45\,826$.
}
\end{tcolorbox}
\end{minipage}
\hfill
\begin{minipage}[t]{0.49\textwidth}\vspace{0pt}
\begin{tcolorbox}[
  enhanced,
  colback=red!4!white,
  colframe=red!70!black,
  arc=2mm, boxrule=0.8pt,
  left=2mm, right=2mm, top=1.8mm, bottom=1.8mm,
  title={\footnotesize\bfseries\color{white}\faExclamationTriangle\ Mẹo tránh bẫy trắc nghiệm}
]
{\footnotesize
\begin{itemize}\setlength{\itemsep}{1.5pt}\setlength{\leftskip}{-2.5mm}
  \item \textbf{Bẫy quên ghi số 0 ở thương:}\\
  $63\,045 : 3$, nếu học sinh quên số 0 sẽ ra $2115$ (giảm đi 10 lần)! Số bị chia có 5 chữ số thì thương thường có 4 hoặc 5 chữ số.
  \item Luôn thử lại bằng phép tính nhân nhẩm hàng cao nhất ($20\,000 \times 3 = 60\,000$).
\end{itemize}
}
\end{tcolorbox}
\end{minipage}


\end{document}
